\documentclass[11pt]{article}
\usepackage{fontspec}
\usepackage[french]{babel}
\usepackage{amsmath}
\usepackage{amssymb}
\usepackage{amsthm}
\usepackage{mathrsfs}
\usepackage{tikz-cd}
\usepackage{enumitem}

\newcommand{\printedpage}[1]{\marginpar{\footnotesize p.~#1}}
\newcommand{\msp}[1]{\marginpar{\footnotesize [#1]}}
\newcommand{\uHom}{\underline{\mathrm{Hom}}}
\newcommand{\uFl}{\underline{\mathrm{Fl}}}
\newcommand{\uIs}{\underline{\mathrm{Is}}}
\newcommand{\uGamma}{\underline{\Gamma}}
\newcommand{\bDelta}{\boldsymbol{\Delta}}
\newcommand{\Ens}{\mathrm{Ens}}
\newcommand{\Ob}{\mathrm{Ob}\,}
\newcommand{\Fl}{\mathrm{Fl}}
\newcommand{\Hom}{\mathrm{Hom}}

\begin{document}

\printedpage{154}
% head: A. Grothendieck : Dérivateurs, Ch. XVIII, Édition des Universités de Montpellier II et Paris VII
dans la catégorie des double-flèches de $\mathcal{E}_s$ ($I$ désignant la catégorie double-flèche type $0\rightrightarrows 1$), foncteur qui est (trivialement) accessible. Soit alors $I\xrightarrow{\;i\;}\bDelta_1$ le foncteur canonique, et
\[
\underbrace{\uHom(\bDelta_1,\mathcal{E}_s)}_{\simeq\,\uFl(\mathcal{E}_s)}\ \hookrightarrow\ \uHom(I,\mathcal{E}_s)
\]
le foncteur $u\longmapsto u\circ i$, qui est un foncteur strictement pleine fidèle, ayant comme image la sous-catégorie des double-flèches $(u,v)$ telles que $u=v$. Comme le foncteur d'inclusion est accessible et $\uHom(\bDelta_1,\mathcal{E}_s)$ et $\uHom(I,\mathcal{E}_s)$ accessibles
\msp{page 406}
(cas particulier de 4.14.1, ou aussi de 4.10.8!), on voit que la sous-catégorie strictement pleine de $\uHom(I,\mathcal{E}_s)$ des double-flèches égalisées dans $\mathcal{E}_s$, est accessible dans $\uHom(I,\mathcal{E}_s)$. Donc, son image inverse dans $\mathcal{M}'$ est accessible. Donc l'intersection de ces images inverses pour $s\in\Ob\mathcal{C}$ et $(\lambda,\mu)\in\Fl_2(\mathcal{C})$ variables, l'est également ($\mathcal{C}$ étant petite).

Donc tout revient à prouver que $\mathcal{M}'$ est accessible, ce qui est contenu dans le

% vertical bar printed in the left margin alongside Lemme 4.14.22.1
\textbf{Lemme 4.14.22.1.} \textit{Soient $\Pi$, $\Pi'$ deux catégories accessibles, $F,G:\Pi\rightrightarrows\Pi'$ deux foncteurs accessibles, et $\mathcal{M}'$ la catégorie des couples $(X,u)$, avec $X\in\Ob\Pi$ et $u:F(x)\longrightarrow G(x)$ une flèche de $\Pi'$. Alors $\mathcal{M}'$ est accessible, et le foncteur canonique $\mathcal{M}'\longrightarrow\Pi$ est accessible.} % [sic] (X,u), X \in Ob Pi ... F(x) -> the same letter (X or x) throughout, capital X and lowercase x used for the same object

\textsc{Démonstration du Lemme.} La donnée de $\Pi$, $\Pi'$, $F$, $G$ définit une catégorie cofibrée $\mathcal{F}$ sur la catégorie $I=(0\rightrightarrows 1)$
\msp{page 407}
précédente. La catégorie $\uGamma(\mathcal{F}/I)$ s'interprète, à isomorphisme près, comme la catégorie des triples $(x,u,v)$ avec $x\in\Ob\Pi$, $x'\in\Ob\Pi'$, et $u$ et $v$ des flèches
\[
\begin{tikzcd}[column sep=small]
F(x) \arrow[dr, "u"'] & & G(x) \arrow[dl, "v"] \\
& x' &
\end{tikzcd}
\]
dans $\Pi'$. La catégorie $\mathcal{M}'$ est équivalente à la sous-catégorie strictement pleine de $\uGamma(\mathcal{F}/I)$, formée des triples $(x,u,v)$ tels que $v$ soit un isomorphisme, i.e. à la sous-catégorie de $\uGamma(\mathcal{F}/I)$, image inverse de $\uIs(\Pi')$ par le foncteur $(x,u,v)\longmapsto v$,
\[
\uGamma(\mathcal{F}/I)\ \longrightarrow\ \uFl(\Pi')\ .
\]
Or par 4.10.8, $\uGamma(\mathcal{F}/I)$ est accessible, et le foncteur précédent est bien sûr accessible. Enfin, par 4.14.12,
\msp{page 408}
$\uIs(\Pi')$ est accessible dans la catégorie accessible $\uFl(\Pi')$. Donc son image inverse $\mathcal{M}'$ dans $\uGamma(\mathcal{F}/I)$ est accessible dans $\uGamma(\mathcal{F}/I)$, donc accessible. Donc il en est de même de la catégorie équivalente, q.e.d.

Ainsi, grâce aux résultats généraux sur les parties accessibles des catégories accessibles, le théorème 4.14.17, dans le cas où $p$ est fibrant, est ramené aux cas cofibrant, dans le cas très particulier où la catégorie d'indices est $I=(0\rightrightarrows 1)$. On pourrait d'ailleurs faire une réduction similaire dans le cas où $p$ est cofibrant, par une démonstration pratiquement % [sic] aux cas cofibrant -> au cas cofibrant, singular noun
% foot: Version 17 décembre 2025
\printedpage{155}
% head: A. Grothendieck : Dérivateurs, Ch. XVIII, Édition des Universités de Montpellier II et Paris VII
identique. En d'autres termes, le théorème 4.10.8 se ramène au cas où la catégorie d'indices est $I=(0\rightrightarrows 1)$. Mais je suspecte que la démonstration du théorème dans ce cas particulier, n'est pas essentiellement plus simple que dans le cas général.

\msp{page 409}

\textbf{4.14.23.} On peut se demander de trouver un énoncé assurant l'accessibilité d'une catégorie $\uGamma(\mathcal{E}/\mathcal{C})$, qui coiffe les deux cas envisagés dans 4.14.17. Voici ce que je peux proposer. Je vais faire l'hypothèse suivante sur les foncteur $p:\mathcal{E}\longrightarrow\mathcal{C}$ (en plus de celle que $\mathcal{C}$ soit essentiellement petite) : % [sic] les foncteur -> le foncteur, singular functor p

\begin{itemize}[leftmargin=*]
\item[H)] Pour tout $\lambda\in\Fl\mathcal{C}$, la catégorie $\mathcal{E}_\lambda=\mathcal{E}\times_{\mathcal{C}}(\bDelta_1,i_\lambda)$, où $i_\lambda:\bDelta_1\longrightarrow\mathcal{C}$ est le foncteur défini par $\lambda$, est accessible, et le foncteur canonique $\mathcal{E}_\lambda\xrightarrow{\;p_\lambda\;}\bDelta_1$ est accessible.
\end{itemize}

Notons que si $\lambda:s\longrightarrow t$ dans $\mathcal{C}$ est tel que $\mathcal{E}_\lambda$ et $p_\lambda$ soient accessibles, alors $\mathcal{E}_s$ et $\mathcal{E}_t$ sont accessible, car isomorphes aux catégories image inverse par $p_\lambda:\mathcal{E}_\lambda\longrightarrow\bDelta_1$ des sous-catégories $\{0\}$ resp. $\{1\}$ de $\Delta_1$, lesquelles sont accessibles dans la catégorie accessible $\bDelta_1$, donc % [sic] sont accessible -> sont accessibles, plural agreement
\msp{page 410}
$\mathcal{E}_{\lambda_0}$ et $\mathcal{E}_{\lambda_1}$ sont accessibles dans $\mathcal{E}_\lambda$, donc sont accessibles, donc aussi $\mathcal{E}_s$ et $\mathcal{E}_t$. Ainsi, l'hypothèse H) implique que les catégories fibres sont accessibles. Notons de plus :

\textbf{Lemme 4.14.23.1.} \textit{Supposons que $p:\mathcal{E}\longrightarrow\mathcal{C}$ soit cofibrant ou fibrant, et soit $(\lambda:s\longrightarrow t)\in\Fl(\mathcal{C})$. Pour que $\mathcal{E}_\lambda$ soit accessible et $\mathcal{E}_\lambda\xrightarrow{\;p_\lambda\;}\bDelta_1$ accessible, il faut et il suffit que $\mathcal{E}_s$, $\mathcal{E}_t$ soient accessibles et $\lambda_*:\mathcal{E}_s\longrightarrow\mathcal{E}_t$ resp. $\lambda^*:\mathcal{E}_t\longrightarrow\mathcal{E}_s$ soit un foncteur accessible.}

\textsc{Démonstration.} On peut supposer que $\mathcal{C}=\bDelta_1$, et que $\mathcal{E}$ est la catégorie cofibrée (resp. fibrée) sur $\bDelta_1$ définie par un foncteur $\varphi:\mathcal{E}_0\longrightarrow\mathcal{E}_1$ (resp. $\psi:\mathcal{E}_1\longrightarrow\mathcal{E}_0$), et à montrer que $\mathcal{E}$ et $p:\mathcal{E}\longrightarrow\bDelta_1$ sont accessibles si et seulement si $\mathcal{E}_0$ et $\mathcal{E}_1$ et $\varphi$ (resp. $\psi$) le sont.

\msp{page 411}

Prouvons le “il faut”. On a déjà vu que si $\mathcal{E}$, $p$ sont accessibles, $\mathcal{E}_0$ et $\mathcal{E}_1$ sont accessibles dans $\mathcal{E}_s$, donc accessibles, et les foncteurs d'inclusion % [sic] dans E_s -> dans E, here C = Delta_1 and the subcategories lie in E
\[
\mathcal{E}_0\ \hookrightarrow\ \mathcal{E}\ \hookleftarrow\ \mathcal{E}_1
\]
accessibles. Dire que $p$ est cofibrant (resp. fibrant) signifie que le foncteur d'inclusion $\mathcal{E}_1\hookrightarrow\mathcal{E}$ admet un adjoint à droite $\bar{\varphi}$ (resp. que le foncteur $\mathcal{E}_0\hookrightarrow\mathcal{E}$ admet un adjoint à gauche $\tilde{\psi}$). Alors $\varphi$ (resp. $\psi$) est le composé $\mathcal{E}_0\longrightarrow\mathcal{E}\xrightarrow{\;\bar{\varphi}\;}\mathcal{E}_1$ (resp. le composé $\mathcal{E}_1\hookrightarrow\mathcal{E}\xrightarrow{\;\tilde{\psi}\;}\mathcal{E}_1$). Donc pour prouver que $\varphi$ (resp. $\psi$) est accessible, il suffit de prouver que $\bar{\varphi}$ (resp. $\tilde{\psi}$) l'est. Cela est un cas particulier du % [sic] E -> E_1 (target of psi-tilde composite) -> E_0, psi goes from E_1 to E_0

\textbf{Lemme 4.14.23.2.} \textit{Soient $\mathcal{E}$, $\mathcal{E}'$ deux catégories accessibles} (\footnote[172]{\textbf{N.B.} Il suffit que $\mathcal{E}$ admette une petite sous-catégorie génératrice, et que les objets de $\mathcal{E}'$ soient accessibles.}), \textit{et $\mathcal{E}\underset{G}{\overset{F}{\rightleftarrows}}\mathcal{E}'$ un couple de foncteurs adjoints entre $\mathcal{E}$ et $\mathcal{E}'$. Alors $F$ et $G$ sont accessibles.}

\msp{page 412}

En effet, $F$ est accessible, car il commute aux petites $\varinjlim$ quelconques, a fortiori il satisfait $\mathrm{L}_{\pi_0}$ (où $\pi_0$ est tel que $\mathcal{E}$, $\mathcal{E}'$ satisfassent $\mathrm{L}_{\pi_0}$). Prouvons que $G$ est accessible. Si $S$ est un
% foot: Version 17 décembre 2025
\printedpage{156}
% head: A. Grothendieck : Dérivateurs, Ch. XVIII, Édition des Universités de Montpellier II et Paris VII
petit ensemble d'éléments générateur dans $\mathcal{E}$, de sorte que la petite famille des foncteurs
\[
h_\xi\ =\ (x\longmapsto\Hom_{\mathcal{E}}(\xi,x))\ :\ \mathcal{E}\ \longrightarrow\ \Ens
\]
est conservative, pour prouver que $G$ est accessible, il suffit de prouver que les $h_\xi\circ G$ ($\xi\in S\subset\Ob\mathcal{E}$) le sont. Or
\[
L_\xi\circ G(x')\ =\ \Hom(\xi,G(x'))\ \simeq\ \Hom_{\mathcal{E}'}(F(\xi),x')\ , % [sic] L_xi -> h_xi, the functors were named h_xi
\]
et comme $F(\xi)$ est un objet accessible de $\mathcal{E}'$, le foncteur
\[
x'\ \longmapsto\ (L_\xi\circ G)(x')\ \simeq\ \Hom_{\mathcal{E}'}(F(\xi),x') % [sic] L_xi -> h_xi, the functors were named h_xi
\]
est bien accessible, q.e.d.

Cela prouve le “il faut” dans 4.14.23.1. Prouvons le “il suffit”, i.e. supposons $\mathcal{E}_s$, $\mathcal{E}_t$ accessibles, et $\varphi$ (resp. $\psi$) accessible.

\msp{page 413}

Dans les deux cas, $\mathcal{E}_0$ est un ouvert de $\mathcal{E}$, $\mathcal{E}_1$ est le fermé complémentaire, et $\mathcal{E}$ est la catégorie “recollée” à l'aide du bifoncteur
\[
H\ :\ \mathcal{E}_0^{\circ}\times\mathcal{E}_1\ \longrightarrow\ \Ens
\]
donné par
\[
\left\{
\begin{aligned}
H(x_0,x_1)\ &=\ \Hom(\varphi(x_0),x_1)\\
&\hspace{-2em}\text{resp.}\\
H(x_0,x_1)\ &=\ \Hom(x_0,\psi(x_1))\ ,
\end{aligned}
\right.
\]
suivant qu'on est dans le cas $p$ cofibrant ou $p$ fibrant. On va utiliser la

\textbf{Proposition 4.14.24} (\footnote[173]{cf. énoncés plus généraux et démonstration sympathique, 4.15.1 et 4.15.4 (p. 425, 432).}). \textit{Soit $\mathcal{E}$ une catégorie, $\mathcal{E}_0$ une sous-catégorie ouverte, $\mathcal{E}_1$ la sous-catégorie fermée complémentaire. Conditions équivalentes :}
\begin{itemize}[leftmargin=*]
\item[\textit{a)}] \textit{$\mathcal{E}$ est accessible, et $\mathcal{E}\longrightarrow\bDelta_1$ (de fibres $\mathcal{E}_0$, $\mathcal{E}_1$) est accessible.}
\item[\textit{b)}] \textit{$\mathcal{E}$ est accessible, et $\mathcal{E}_0$, $\mathcal{E}_1$ sont accessibles dans $\mathcal{E}$.}
\item[\textit{c)}] \textit{$\mathcal{E}_0$, $\mathcal{E}_1$ sont accessibles, et pour tout $x_0\in\Ob\mathcal{E}_0$, le foncteur}
\[
x_1\ \longmapsto\ \uHom_{\mathcal{E}}(x_0,x_1)\ :\ \mathcal{E}_1\ \longrightarrow\ \Ens
\]
\msp{page 414}
\textit{est accessible, et il existe $\pi$ tel que pour tout $x_1\in\Ob\mathcal{E}_1$, le foncteur}
\[
x_0\ \longmapsto\ \uHom_{\mathcal{E}}(x_0,x_1)^{\circ}\ :\ \mathcal{E}_0\ \longrightarrow\ (\Ens)^{\circ}
\]
\textit{soit $\pi$-accessible.}
\end{itemize}
% foot: Version 17 décembre 2025

\printedpage{157}
% head: A. Grothendieck : Dérivateurs, Ch. XVIII, Édition des Universités de Montpellier II et Paris VII
Admettons ce lemme \texttt{[plutôt cette proposition]}, achevons la démonstration de 4.14.23.1. On suppose $\mathcal{E}_0$, $\mathcal{E}_1$ et $\varphi$ (resp. $\psi$) accessibles, on veut prouver la condition a) de 4.14.23.3. Pour ceci, il suffit de prouver la condition c). Or on a, dans le cas $p$ cofibrant, % [sic] 4.14.23.3 -> 4.14.24, the conditions a), c) are those of Proposition 4.14.24
\[
\Hom_{\mathcal{E}}(x_0,x_1)\ \simeq\ \Hom_{\mathcal{E}_1}(\varphi(x_0),x_1)\ ,
\]
c'est bien accessible en $x_1$, puisque $\mathcal{E}_1$ est accessible, et $\pi$-accessible en $x_0$ si $\varphi$ l'est, comme composé de $\varphi$ avec $\mathcal{E}_1\longrightarrow\Ens$, $\Hom(-,x_1)$, qui l'est. Dans le cas $p$ fibrant, on a
\[
\Hom_{\mathcal{E}}(x_0,x_1)\ \simeq\ \Hom_{\mathcal{E}_0}(x_0,\psi(x_1))\ ,
\]
c'est accessible en $x_1$, comme le composé du foncteur accessible $\psi:\mathcal{E}_1\longrightarrow\mathcal{E}_0$, et du foncteur également accessible $\Hom(x_0,-):\mathcal{E}_0\longrightarrow\Ens$, qui est $\pi$-accessible en $x_0$ tautologiquement.

Pour achever la démonstration de 4.14.23.1, il reste donc à prouver la proposition 4.14.24. On a déjà noté qu'on a a $\Longrightarrow$ b, et b $\Longrightarrow$ c clair, puisque le foncteur envisagé $\mathcal{E}\longrightarrow\Ens$ est le composé $\mathcal{E}_1\overset{\text{incl.}}{\hookrightarrow}\mathcal{E}\xrightarrow{\;\Hom_{\mathcal{E}}(x_0,-)\;}\Ens$, est accessible comme composé de deux
\msp{page 415}
foncteurs accessibles. Reste donc à prouver c $\Longrightarrow$ a. Soit donc $\pi_0$ un cardinal utile pour $\mathcal{E}_0$, et tel que les foncteurs
\[
x_0\ \longmapsto\ H(x_0,x_1)\ :\ \mathcal{E}_0\longrightarrow\Ens^{\circ}
\]
soient $\pi_0$-accessibles, soit $\pi$ un cardinal $\geqslant\sup(\pi_0,\mathrm{card}\,\text{réd}\,\mathcal{E}_{0\,\pi_0})$ tel que $\pi^{\pi_0}=\pi$ (donc utile pour $\mathcal{E}_0$), utile pour $\mathcal{E}_1$, et tel que les foncteurs
\[
x_1\ \longmapsto\ H(x_0,x_1)\ :\ \mathcal{E}_1\longrightarrow\Ens
\]
soient $\pi$-accessibles, pour $x_0\in\Ob\mathcal{E}_{0\,\pi_0}$. Soit $\mathcal{E}(\pi)$ la sous-catégorie strictement pleine de $\mathcal{E}$, telle que
\[
\Ob\mathcal{E}(\pi)\ =\ \Ob\mathcal{E}_{0\,\pi}\cup\Ob\mathcal{E}_{1\,\pi}\ .
\]
Je dis que
\begin{itemize}[leftmargin=*]
\item[1°)] $\mathcal{E}$ satisfait à $\mathrm{L}_\pi$, et $\mathcal{E}_0$ et $\mathcal{E}_1$ sont stables dans $\mathcal{E}$ sous $\mathrm{L}_\pi$, les foncteurs $\mathcal{E}_0\longrightarrow\mathcal{E}$, $\mathcal{E}_1\longrightarrow\mathcal{E}$, $\mathcal{E}\longrightarrow\bDelta_1$ satisfont $\mathrm{L}_\pi$ (\footnote[174]{il suffit que $\mathcal{E}_0$ et $\mathcal{E}_1$ satisfassent $\mathrm{L}_\pi$, et que
\[
x_0\ \longmapsto\ H(x_0,x_1)\ :\ \mathcal{E}_0\longrightarrow\Ens^{\circ}
\]
satisfasse $\mathrm{L}_\pi$ $\forall\ x\in\Ob\mathcal{E}$.}).
\item[2°)] $\mathcal{E}(\pi)=\mathcal{E}_\pi$ (d'où un foncteur pleinement fidèle
\[
\mathrm{Ind}_\pi(\mathcal{E}(\pi))\ \longrightarrow\ \mathcal{E}\ \ ).
\]
\item[3°)] Ledit foncteur est essentiellement surjectif,
\msp{page 416}
donc une équivalence de catégories.
\end{itemize}
Notons que 2°) et 3°) impliquent que $\mathcal{E}$ est accessible, et 1°) que $\mathcal{E}\longrightarrow\bDelta_1$ est accessible, d'où la condition a) du lemme qu'il s'agit d'établir. Donc tout revient à prouver les 1°, 2°, 3° ci-dessus.
% foot: Version 17 décembre 2025

\printedpage{158}
% head: A. Grothendieck : Dérivateurs, Ch. XVIII, Édition des Universités de Montpellier II et Paris VII
1°) Comme $\mathcal{E}_1$ est fermée dans $\mathcal{E}$, il est clair que $\mathcal{E}$ est stable par toutes $\varinjlim$ représentables dans $\mathcal{E}$, et plus précisément, que pour un foncteur $I\longrightarrow\mathcal{E}_1$, sa limite inductive dans $\mathcal{E}_1$ existe si et seulement si elle existe dans $\mathcal{E}$, et les deux sont isomorphes. Supposons $I$ grande devant $\pi$, et soit $\varphi:I\longrightarrow\mathcal{E}$, $i\longmapsto x_i$. Distinguons deux cas.

a) $\varphi$ prend ses valeurs dans $\mathcal{E}_0$. Soit $x_0$ sa limite inductive \textit{dans} $\mathcal{E}_0$. Je dis que c'est aussi une limite inductive dans $\mathcal{E}$. Il faut prouver que pour $x_1\in\Ob\mathcal{E}_1$, la flèche
\[
\Hom_{\mathcal{E}}(x_0,x_1)\ \longrightarrow\ \varprojlim_{i\in I}\Hom_{\mathcal{E}}(x_i,x_1)
\]
est isomorphe. Or c'est le cas
\msp{page 417}
par hypothèse sur $\pi$,
\[
\xi\ \longmapsto H(\xi,x_1)\ :\ \mathcal{E}_0\ \longrightarrow\ (\Ens)^{\circ}
\]
satisfait à $\mathrm{L}_\pi$ pour tout $x_1$ dans $\Ob\mathcal{E}_1$.

b) $\varphi$ ne se factorise pas par $\mathcal{E}_0$, i.e. $\exists\ i_0\in\Ob I$ avec $\varphi(i_0)\in\Ob\mathcal{E}_1$. Alors on a $\varphi(i)\in\Ob\mathcal{E}_1$ pour tout $i\in\Ob I[i_0]$, où $I[i_0]$ désigne la sous-catégorie strictement pleine de $I$ formée des $i$ tels que $i\geqslant i_0$. Cette catégorie est cofinale dans $I$, donc grande devant $\pi$ comme $I$, et pour que $\varinjlim\limits_{I}x_i$ existe, il faut et il suffit que $\varinjlim\limits_{I[i_0]}x_i$ existe, ce qui est le cas d'après les cas a) déjà traité. % [sic] les cas a) déjà traité -> le cas a) déjà traité, singular agreement

Cela prouve donc en même temps que $\mathcal{E}$ satisfait à $\mathrm{L}_\pi$, et que $\mathcal{E}_0$ et $\mathcal{E}_1$ sont $\mathrm{L}_\pi$-stables. La démonstration montre en même temps que le foncteur $p:\mathcal{E}\longrightarrow\bDelta_1$ satisfait à $\mathrm{L}_\pi$.

2°) De ceci résulte aussitôt que $\mathcal{E}_\pi\subset\mathcal{E}(\pi)$. Prouvons qu'on a aussi
\[
\mathcal{E}(\pi)\ \subset\ \mathcal{E}_\pi\ ,
\]
et distinguons encore deux cas, pour
\msp{page 418}
$x\in\Ob E(\pi)$, suivant que $x$ est dans $\mathcal{E}_{0\,\pi}$ ou $\mathcal{E}_{1\,\pi}$. S'il est dans $\mathcal{E}_{1\,\pi}$, pour prouver que % [sic] E(pi) -> \mathcal{E}(pi), script letter elsewhere
\[
\Hom_{\mathcal{E}}(x,\underbrace{\varinjlim y_i}_{y})\ \xleftarrow{\ \sim\ }\ \varinjlim_{i}\Hom_{\mathcal{E}}(x,y_i)
\tag*{$(*)$}
\]
(pour un système inductif $(y_i)_{i\in I}$ dans $\mathcal{E}$ grand devant $\pi$), si les $y_i$ sont dans $\mathcal{E}_0$, donc $y\in\mathcal{E}_0$, alors les deux membres sont vides, et c'est OK. Sinon, on peut supposer (quitte à remplacer $I$ par un $I[i_0]$) que les $x_i$ sont dans $\mathcal{E}_1$, et alors c'est OK aussi, puisque $y$ est aussi la $\varinjlim$ dans $\mathcal{E}_1$, d'après 1°). % [sic] les x_i -> les y_i, the inductive system is (y_i)

Si $x$ est dans $\mathcal{E}_{0\,\pi}$, alors $(*)$ est clair si les $y_i$ sont dans $\mathcal{E}_0$. Sinon, on peut supposer encore que les $y_i$ sont dans $\mathcal{E}_0$, et $(*)$ s'écrit % [sic] y_i sont dans E_0 (second occurrence) -> E_1, the case "sinon" and the following "les y_i dans E_1"
\[
H(x,\varinjlim y_i)\ \xleftarrow{\ \sim\ }\ \varinjlim H(x,y_i)\ ,
\]
pour $x\in\Ob\mathcal{E}_{0\,\pi}$, les $y_i$ dans $\mathcal{E}_1$, $I$ grande devant $\pi$. En d'autres termes, il faut prouver que le foncteur
\[
x_1\ \longmapsto\ H(x_0,x_1)\ :\ \mathcal{E}_1\ \longrightarrow\ \Ens
\]
% foot: Version 17 décembre 2025

\end{document}
